% Keep the interpretation on its prescribed fifth paper page.
\clearpage
\section{Analysis and Discussion}
\label{sec:analysis}

The staged results do not form a monotonic accuracy ladder. Instead, they show
which aspects of the prediction task changed as the formulation progressed
from individual insertion-loss values to complete complex matrices.

\subsection{Point-Wise Capacity Was Not Sufficient}
The first finding is that generic point-wise complexity did not consistently
reduce selected-path error in the retained configurations. Scalar and Vector
Ridge necessarily gave the same $IL_{7,1}$ solution, while the polynomial,
forest and MLP variants remained close to, or worse than, the stronger
frequency-dependent mean-curve reference. This suggests that adding capacity
without a more appropriate broadband representation did not reliably extract
additional design-conditioned information for this path. It does not show that
nonlinear point-wise models are generally ineffective, because only selected
configurations and one split were evaluated.

Schierholz \textit{et al.} demonstrated direct prediction of selected
S-parameter magnitude and phase on the source database, while reporting greater
difficulty for electrically long, strongly frequency-dependent links
\cite{schierholz2021sipi}. The present sharp-null errors are consistent with
that qualitative difficulty, although the different targets and protocols
preclude a numerical ranking. Similarly, Swaminathan \textit{et al.} described
the poor scaling of conventional fully connected networks for broadband SI/PI
responses \cite{swaminathan2020demystifying}. The narrow spread among the
point-wise models here provides project-specific motivation for changing the
data interface rather than only enlarging the regressor.

\subsection{Whole-Curve Modelling Improved the Selected Path}
The whole-curve model produced the lowest selected $IL_{7,1}$ error, and its
paired improvement over Polynomial MLP crossed zero but not the complete
predeclared practical-equivalence band. Joint prediction therefore gave the
strongest retained selected-path result, but the evidence supports association,
not a causal claim about frequency structure: the interface, scaling,
architecture and loss changed together.

This result is directionally consistent with the S-TCNN principle of Torun
\textit{et al.}, which uses transposed convolution to exploit correlation along
a generated spectral response \cite{torun2019spectral}. The current decoder
adapts that principle to six PCB insertion-loss curves rather than reproducing
their architecture or application. Consequently, the comparison supports the
design rationale but cannot establish superiority across datasets. The
remaining deep-null errors also indicate that whole-curve structure alone did
not resolve weak-response regions.

\subsection{Scope, Accuracy and Physical Validity Are Separate}
The full-S-matrix model answered a broader question: it learned
design-conditioned reflection, transmission, crosstalk and phase information
for every port pair, improved on a same-task mean matrix, and guaranteed
reciprocity. Its higher extracted $IL_{7,1}$ error does not prove a cost of
larger output scope, because the curve and matrix models also used different
architectures, scaling and objectives. Selected-path accuracy and response
coverage must therefore be reported as distinct axes.

Physical validity forms a third axis. Chen and Tan showed that low numerical
error can coexist with non-physical insertion-loss responses and therefore
motivated physics-enforced evaluation \cite{chen2021physics}. The absence of
sampled passivity violations here is encouraging but weaker evidence: no
passivity term was trained and behaviour outside the evaluated grid is unknown.
Akinwande \textit{et al.} predicted selected complex broadband responses with a
complex-valued network and explicit passivity and causality treatment
\cite{akinwande2024complex}. In contrast, the present real-valued model covers
a complete 12-port matrix but enforces only reciprocity. Its larger predicted
than true finite-band causality residual reinforces that architectural
reciprocity cannot substitute for passivity or causality constraints.

\subsection{Confidence and Scope}
The evidence is conditional on one selected training seed, so optimisation
variability is unknown. Test designs were excluded from fitting, but their
results were inspected during development; the bootstrap intervals quantify
variation across those fixed test designs, not uncertainty from retraining,
model search or repeated exposure. Random held-out designs from one topology
and sampled parameter ranges do not establish extrapolation, cross-topology
generalisation or task-specific accuracy sufficient for an SI decision. The
passivity result covers only sampled matrices, and the positive finite band
limits the causality diagnostic. Finally, the planned controlled parameter
study and electromagnetic-simulation timing comparison were not completed, so
neither physical sensitivity nor computational saving can be inferred.
